% Addendum to rigid_subcomplexes.tex (2026-10-07): rigid graphs over minimal complexes, maximal star
% rectangles, and the exact completion of the icosahedron stall.  Uses the macros and theorem environments
% of rigid_subcomplexes.tex (\SCst, \TC, \gen, theorem/lemma/proposition/corollary/remark).

\section{Rigid graphs over minimal complexes}

\begin{definition}
A vertex $v$ of $K$ is \emph{dominated} if some vertex $v'\ne v$ belongs to every facet of $K$ that contains
$v$.  Thus $K$ has no dominated vertex if and only if $\bigcap\{\sigma\ \text{facet}: v\in\sigma\}=\{v\}$ for
every vertex $v$.
\end{definition}

Every closed surface triangulation has no dominated vertex (the facets at $v$ are the triangles $\{v,a_i,a_{i+1}\}$
around the link cycle of $v$, of length at least~$3$, and their intersection is $\{v\}$); so does
$\partial\Delta^{N-1}$ and the five-vertex M\"obius band.  More generally a complex without free faces has no
dominated vertex, since removing a dominated vertex is a strong collapse, hence a sequence of elementary collapses.

\begin{lemma}\label{lem:idalone}
If $K$ has no dominated vertex and $\varphi:K\to K$ is a simplicial map contiguous to the identity, then
$\varphi=\mathrm{id}$.
\end{lemma}
\begin{proof}
For a facet $\sigma$, contiguity gives $\sigma\cup\varphi(\sigma)\in K$; since $\sigma$ is maximal,
$\varphi(\sigma)\subseteq\sigma$.  Hence for a vertex $v$,
$\varphi(v)\in\bigcap_{\sigma\ni v}\sigma=\{v\}$.
\end{proof}

\begin{theorem}[rigid graphs]\label{thm:graphs}
Let $K$ have no dominated vertex and let $w:V\to V$ be non-decreasing for the vertex order and simplicial
($w(\sigma)\in K$ for every $\sigma\in K$), with $w\ne\mathrm{id}$.  Put
\[
\Gamma_w=\bigl\{\{(v,w(v)):v\in\sigma\}:\sigma\in K\bigr\}.
\]
Then $\Gamma_w$ is a subcomplex of $K\times K$, $\pi_1|_{\Gamma_w}$ is an isomorphism onto $K$ under which
$\pi_2|_{\Gamma_w}$ becomes $w$, and $\Gamma_w$ is rigid.  The same holds for the transpose
$\Gamma_w^{\mathsf T}=\{\{(w(v),v):v\in\sigma\}:\sigma\in K\}$.
\end{theorem}
\begin{proof}
For $\sigma\in K$ let $S_\sigma=\{(v,w(v)):v\in\sigma\}$.  Since $w$ is non-decreasing, $S_\sigma$ is a chain in
$V\times V$; $\pi_1(S_\sigma)=\sigma\in K$ and $\pi_2(S_\sigma)=w(\sigma)\in K$ because $w$ is simplicial.  By
Lemma~\ref{lem:criterion} $S_\sigma$ is a simplex of $K\times K$, and $S_\tau\subseteq S_\sigma$ for
$\tau\subseteq\sigma$, so $\Gamma_w$ is a subcomplex.  The map $(v,w(v))\mapsto v$ is a bijection on vertices
carrying $S_\sigma$ to $\sigma$, hence an isomorphism $\Gamma_w\to K$; under it $\pi_1$ becomes $\mathrm{id}_K$ and
$\pi_2$ becomes $w$.  A contiguity chain $\pi_1|_{\Gamma_w}=\psi_0\sim\dots\sim\psi_L=\pi_2|_{\Gamma_w}$ therefore
becomes a chain $\mathrm{id}=\varphi_0\sim\varphi_1\sim\dots\sim\varphi_L=w$ of simplicial maps $K\to K$.  By
Lemma~\ref{lem:idalone} and induction, $\varphi_i=\mathrm{id}$ for all $i$, so $w=\mathrm{id}$, a contradiction.
For the transpose read the chain backwards: it starts at the isomorphism $\pi_2|_{\Gamma_w^{\mathsf T}}$.
\end{proof}

\begin{remark}
(a) For $K=\partial\Delta^{N-1}$ every non-decreasing $w\ne\mathrm{id}$ is simplicial (it is not onto by
Lemma~\ref{lem:notonto}, so proper subsets go to proper subsets) and $\Gamma_w=Z_w$: Theorem~\ref{thm:graphs}
contains Proposition~\ref{prop:spheres} and extends it to every closed surface triangulation.
(b) Constant $w\equiv c$ gives $\Gamma_c=K\times\{c\}$: no contiguity domain contains a full row or column of
$K\times K$.
(c) If $K$ is contractible and has no free face (a dunce-hat triangulation), then $K\times K\supseteq K\times\{c\}$
is not a domain, so $\SCst(K)\ge1>0=\TC(|K|)$ although $\pi_1$ and $\pi_2$ are homotopic on $K\times\{c\}$: the
rigidity of Theorem~\ref{thm:graphs} is combinatorial, not homotopical.
\end{remark}

\section{Star rectangles are maximal domains}

Write $D(v,w)$ for the star rectangle of the preprint and recall that a rigid subcomplex of $\gen{D(v,w)\cup\{f\}}$
shows that no domain contains $D(v,w)\cup\{f\}$ (Lemma~\ref{lem:hered}).

\begin{theorem}[computer-verified]\label{thm:maximal}
Let $K$ be the hemi-icosahedron, the five-vertex M\"obius band $K_M$, or the seven-vertex (Cs\'asz\'ar) torus.
For every vertex pair $(v,w)$ and every facet $f$ of $K\times K$ not in $D(v,w)$, the subcomplex
$\gen{D(v,w)\cup\{f\}}$ contains a hollow triangle $Z\cong C_3$ whose two projected closed walks $\pi_1\circ Z$,
$\pi_2\circ Z$ lie in different contiguity classes of maps $C_3\to K$.  Hence no contiguity domain properly
contains a star rectangle: every $D(v,w)$ is a maximal contiguity domain.
\end{theorem}
\begin{proof}[Proof (computer-assisted)]
Rigidity of such a $Z$ is Lemma~\ref{lem:hered} applied to the restriction: a chain on $Z$ from $\pi_1|_Z$ to
$\pi_2|_Z$ would be a contiguity chain between the two walks.  The $36\cdot450$, $25\cdot96$ and $49\cdot960$
witnesses were produced by \texttt{scripts/rigid\_analysis/rigid.py} and re-verified independently: each edge of
$Z$ is checked to lie in a facet of $D(v,w)\cup\{f\}$, and the contiguity classes of all maps $C_3\to K$ are
recomputed by brute force from Definition~\ref{def:contig} ($2$, $3$ and $7$ classes respectively).  The
certificates are \texttt{scripts/rigid\_analysis/certificates/maximal\_rectangles\_\{rp2,mobius,torus\}.json}.
Every witness is also a hand proof via Lemma~\ref{lem:winding}: for the hemi-icosahedron, $D(1,1)$ and the facet
over $(\{0,4,5\},\{2,3,5\})$, the points $(0,0),(0,2),(5,5)$ span a hollow triangle with $\pi_1$-walk $0,0,5$
(winding $0$) and $\pi_2$-walk $0,2,5$, the essential empty triangle $\{0,2,5\}$ (winding $1$).
\end{proof}

\begin{corollary}
On these three complexes a system of $k$ chains whose domains are star rectangles cannot increase its coverage by
changing a single chain, whatever the change: any domain containing a rectangle equals it.  In particular the
annealer's plateau at $k=\rho-1$ on the hemi-icosahedron ($540/600$, five diagonal rectangles) is a strict local
maximum of coverage under all single-chain moves.
\end{corollary}

\section{The icosahedron: a stall that was not an obstruction}

\begin{proposition}\label{prop:icosa}
$\SCst(\text{icosahedron boundary})\le3$.
\end{proposition}
\begin{proof}
Certificate \texttt{results/icosa\_k4\_exact}: four domains of $819$, $728$, $629$ and $224$ facets with chains of
$8,9,7,6$ links, verified link by link by \texttt{src/independent\_check.py}.  It was obtained from the stalled
state \texttt{icosa\_k4\_v3} ($2376/2400$ facets): its domain of $744$ facets together with all $24$ uncovered
facets admits a chain of $8$ links (an exact SAT decision, \texttt{scripts/rigid\_analysis/domain\_chain\_sat.py}),
while its smallest domain plus the same $24$ facets admits none of $\le12$ links.
\end{proof}

\begin{remark}
The icosahedron has no empty triangle, all closed walks of length $\le5$ in it are in one contiguity class, and the
family $\Gamma_w$ of Theorem~\ref{thm:graphs} ($3891$ maps) is avoidable by two facet classes; so no rigid
subcomplex known to us obstructs even two domains, and the stall at four was a property of the search, not of the
complex.  Whether three domains suffice ($\TC(S^2)=2$) is being decided by the exact cover search
\texttt{scripts/rigid\_analysis/cegar\_cover.py}.
\end{remark}
